% 第2章 理论基础
\chapter{理论基础}
\label{chap:theory}

本章系统阐述阿秒电子显微术所涉及的核心物理原理，包括电子波函数的光场调制机制、时间Talbot效应与阿秒脉冲形成、高次谐波产生原理、以及电子能量损失谱在近场成像中的应用。

\section{电子波函数的光场调制}
\label{sec:electron-modulation}

\subsection{自由电子在光场中的有质动力相互作用}
\label{subsec:ponderomotive}

当电子束穿越电磁场时，电子波函数会受到场的空间和时间周期性调制。在有质动力（ponderomotive）近似下，电子感受到的等效势能与光场强度的空间梯度成正比：

\begin{equation}
    \label{eq:ponderomotive}
    U_p(\mathbf{r}, t) = \frac{e^2}{4 m_e \omega^2} |\mathbf{E}(\mathbf{r}, t)|^2
\end{equation}

其中$e$为电子电荷，$m_e$为电子质量，$\omega$为光场角频率，$\mathbf{E}$为电场强度。

当电子以速度$v_z$穿越驻波场时，在电子的共动坐标系中，空间周期性转化为时间周期性，电子波函数获得与光子能量$\hbar\omega$对应的能量侧带。量子力学上，这一过程可以用电子与光子间的能量交换描述：

\begin{equation}
    \label{eq:sideband}
    |\psi\rangle = \sum_{n=-\infty}^{\infty} J_n(2|g|) \, |E_0 + n\hbar\omega\rangle
\end{equation}

其中$J_n$为$n$阶贝塞尔函数，$g$为耦合强度参数，$E_0$为电子初始动能，$n\hbar\omega$为第$n$阶侧带的能量偏移。耦合强度$|g|$与光场强度和相互作用长度成正比。

\subsection{能量调制的演化}
\label{subsec:energy-modulation}

初始的正弦能量调制在自由空间传播过程中经历色散演化。在近轴近似下，电子波包的纵向相移与传播距离$z$的关系为：

\begin{equation}
    \label{eq:phase-evolution}
    \phi_n(z) = \frac{n^2 (\hbar\omega)^2}{2 E_0} \cdot \frac{z}{\hbar v_z}
\end{equation}

这一相位演化导致不同阶侧带之间的干涉，使得电子的时间分布（即概率密度沿时间轴的投影）随传播距离发生周期性变化。这正是时间Talbot效应的物理起源。

\section{时间Talbot效应与阿秒脉冲形成}
\label{sec:talbot}

\subsection{经典Talbot效应的回顾}
\label{subsec:classical-talbot}

经典Talbot效应是指：当相干光照射周期性光栅时，在光栅后方特定距离处会出现光栅像的自再现。这一现象最早由Talbot于1836年发现。自再现距离（Talbot距离）为：

\begin{equation}
    \label{eq:talbot-classical}
    z_T = \frac{2 d^2}{\lambda}
\end{equation}

其中$d$为光栅周期，$\lambda$为光波长。

\subsection{时间域Talbot效应}
\label{subsec:temporal-talbot}

Tsarev等人\cite{Tsarev2021}将经典Talbot效应推广到时间域。当电子波函数被激光进行周期性能量调制后，在自由空间传播过程中，不同能量侧带之间的相位演化会导致电子时间分布的自再现。时间Talbot距离为：

\begin{equation}
    \label{eq:talbot-temporal}
    z_T^{\text{time}} = \frac{2 E_0^2}{(\hbar\omega)^2} \cdot \frac{\lambda_e}{2\pi}
\end{equation}

其中$\lambda_e = h / \sqrt{2 m_e E_0}$为电子的德布罗意波长。在$z = z_T^{\text{time}} / 2$处（半Talbot距离），电子的时间分布压缩为原始周期$T = 2\pi/\omega$内的一系列窄脉冲，即阿秒电子脉冲串。图\ref{fig:ch2-talbot}示意了这一过程。

\begin{figure}[htbp]
    \centering
    \includegraphics[width=0.9\linewidth]{ch2/fig_talbot}
    \caption{\label{fig:ch2-talbot} 时间Talbot效应产生阿秒电子脉冲串的原理。(a) 初始能量调制使电子波函数获得周期性侧带；(b) 在半Talbot距离处，各侧带的相长干涉形成阿秒脉冲；(c) 单个脉冲的时间宽度由调制深度和侧带阶数共同决定。}
\end{figure}

\subsection{脉冲宽度的定量估计}
\label{subsec:pulse-duration}

在调制深度$|g|$足够大时（$|g| \gg 1$），阿秒脉冲的时间宽度（FWHM）可近似为：

\begin{equation}
    \label{eq:pulse-width}
    \tau_{\text{as}} \approx \frac{T}{2\pi |g|} = \frac{1}{|g| \cdot \omega}
\end{equation}

对于800\,nm波长的连续波激光（$\omega \approx 2.4 \times 10^{15}\;\text{rad/s}$）和典型调制深度$|g| \approx 3$，脉冲宽度约为$\SI{800}{as}$。增大激光功率可提高$|g|$，从而进一步缩短脉冲宽度。

\section{高次谐波产生的三步模型}
\label{sec:hhg}

\subsection{三步模型的基本图像}
\label{subsec:three-step}

Corkum\cite{Corkum1993}提出的三步模型将高次谐波产生过程分解为以下三个阶段（如图\ref{fig:ch2-hhg}所示）：

\textbf{第一步：隧道电离}。强激光场使原子势垒发生倾斜，价电子通过量子隧穿效应逃离原子实的束缚。电离概率由Landau-Zener公式或ADK（Ammosov-Delone-Krainov）理论描述。

\textbf{第二步：自由电子的加速运动}。电离后的电子在激光场中被加速，获得额外的动能。电子的运动轨迹由激光场的相位决定。根据经典力学，在激光场相位$\phi_0$处电离的电子，其返回时间$t_r$和返回动能$E_{\text{kin}}$满足：

\begin{equation}
    \label{eq:return-energy}
    E_{\text{kin}}(\phi_0) = 2 U_p \left[\sin(\omega t_r - \phi_0) - \sin(-\phi_0)\right]^2
\end{equation}

其中$U_p = e^2 E_0^2 / (4 m_e \omega^2)$为有质动力能量。最大返回动能为$3.17\,U_p$，对应截止频率$\hbar\omega_{\text{cutoff}} = I_p + 3.17\,U_p$（$I_p$为电离能）。

\textbf{第三步：复合辐射}。返回的电子与母离子复合，将获得的动能和电离能以高能光子的形式辐射。辐射光子的能量为$\hbar\omega_n = n\hbar\omega_0$，其中$n$为谐波阶数。

\begin{figure}[htbp]
    \centering
    \includegraphics[width=0.85\linewidth]{ch2/fig_hhg_model}
    \caption{\label{fig:ch2-hhg} 高次谐波产生的三步模型。(a) 原子势在强激光场作用下发生倾斜；(b) 电子隧道电离后被激光场加速；(c) 电子返回母离子并复合辐射出高能光子。}
\end{figure}

\subsection{HHG平台区与阿秒脉冲产生}
\label{subsec:hhg-plateau}

HHG谱具有典型的平台区特征：在低阶谐波区，强度随阶数快速下降；进入平台区后，谐波强度在很宽的能量范围内近似保持恒定；在截止频率处急剧下降。Antoine等人\cite{Antoine1996}指出，平台区内大量等间距谐波的相干叠加可以产生阿秒脉冲串。对于理想的平台区（谐波从$q_{\min}$到$q_{\max}$，等间距，等振幅），合成脉冲的时间宽度为：

\begin{equation}
    \label{eq:atto-train}
    \tau = \frac{2\pi}{(q_{\max} - q_{\min})\omega}
\end{equation}

覆盖的谐波阶数越多，脉冲越短。这一原理是本文理解阿秒时间尺度的基础。

\section{电子能量损失谱与近场成像}
\label{sec:eels}

\subsection{电子能量损失谱的基本原理}
\label{subsec:eels-principle}

高速电子穿过薄样品时，会通过库仑相互作用与样品中的电磁场发生能量交换。电子能量损失谱（electron energy-loss spectroscopy, EELS）记录入射电子在透射过程中损失的能量分布。损失的能量$\Delta E$反映了电子与样品中特定激发模式（等离激元、声子、带间跃迁等）的耦合强度。

\subsection{近场成像的对比度机制}
\label{subsec:nearfield-contrast}

在阿秒电子显微术中，能量滤波器选择获得了特定能量增益$\Delta E > E_{\text{cut}}$的电子。这些电子在与近场相互作用过程中获得了净能量转移。能量增益的空间分布直接反映了纳米结构近场的振幅和相位信息。

与传统的EELS不同，阿秒场周期对比度成像的信号来源是电子与近场中光子态的受激能量交换（photon-induced near-field electron microscopy, PINEM），其能量变化为离散的光子能量$\hbar\omega$的整数倍。通过扫描激光-电子之间的时间延迟$\Delta t$，可以获得近场随时间演化的"电影"。

\begin{figure}[htbp]
    \centering
    \includegraphics[width=0.85\linewidth]{ch2/fig_pinem}
    \caption{\label{fig:ch2-pinem} 光子诱导近场电子显微术（PINEM）的原理。(a) 电子穿过纳米结构的近场区域；(b) 电子与光场交换整数个光子能量，形成能量侧带；(c) 能量滤波后获得近场分布的空间映射。}
\end{figure}

\section{本章小结}
\label{sec:theory-summary}

本章介绍了阿秒电子显微术的理论基础。电子波函数在光场中的有质动力调制产生周期性能量侧带，通过时间Talbot效应在自由空间传播中演化形成阿秒电子脉冲串。高次谐波产生的三步模型为理解阿秒时间尺度提供了物理图像。电子能量损失谱与PINEM机制为近场成像提供了对比度来源。上述理论构成本文实验研究的基础框架。
